% Vallée de stabilité
\begin{minipage}{0.54\linewidth}
Le diagramme ci-dessous donne les noyaux stables des éléments tels que
$1\leqslant Z\leqslant8$.

    \begin{enumerate}
    \item Ces noyaux sont disposés autour de la première bissectrice ($N=Z$).
    \item voir figure ci-contre.
    \item 
          \begin{itemize}
              \item Les noyaux radioactifs $\beta^{-}$ sont en dessus de la
                    bissectrice~: \ch{^{14}C} , \ch{^{16}N}, \ch{^{17}N},
                    \ch{^{19}O}
              \item Les noyaux radioactifs $\beta^{+}$ sont en dessous de la
                    bissectrice~: \ch{^{11}C} , \ch{^{13}N}, \ch{^{15}O}
          \end{itemize}

    \item \ch{^{146}_{6}C -> ^{147}_{6}N + ^{0}_{-1}e}  

           On rejoint la vallée de stabilité en se dirigeant directement vers la
           bissectrice.

    \item \ch{^{12}_{7}N -> ^{12}_{6}N + ^{0}_{+1}e}

          On rejoint la vallée de stabilité par désintégration en se dirigeant
          directement vers la bissectrice.
\end{enumerate}
\end{minipage}
\hfill
\begin{minipage}{0.46\linewidth}
\begin{figure}[H]
    \centering
    \begin{tikzpicture}[]
        \matrix[matrix of nodes,nodes in empty cells,
            nodes={anchor=center,minimum size=6.3mm,draw},
        ] (m)
        {
            &&&&&&&& \\
            &&&&&&&& \\
            &&&&&&&& \\
            &&&&&&&& \\
            &&&&&&&& \\
            &&&&&&&& \\
            &&&&&&&& \\
            &&&&&&&& \\
            &&&&&&&& \\
            &&&&&&&& \\
            &&&&&&&& \\
            &&&&&&&& \\
        };
        \foreach \N [count=\i from 1] in {11,10,...,0}
            \node[left] at (m-\i-1.west) {\N};
        \foreach \Z [count=\i from 1] in {0,...,8}
            \node[below] at (m-12-\i.south) {\Z};
        \draw[Latex-Latex] ([xshift=5mm]m-12-9.south east)
            node[below] {$Z$} -| ([yshift=5mm]m-1-1.north west)
            node[right] {$N$};
        \foreach \i/\j/\X/\A in {2/9/O/18,3/9/O/17,4/9/O/16,4/8/N/15,5/8/N/14,
                                 5/7/C/13,6/7/C/12,6/6/B/11,7/6/B/10,7/5/Be/9,
                                 8/4/Li/7,9/4/Li/6,10/3/He/4,11/3/He/3,
                                 11/2/H/2,12/2/H/1}
            \node[fill=black!15,font=\footnotesize,inner sep=0.6pt,
                minimum size=5.5mm,
                rounded corners=1.4pt] at (m-\i-\j.center) {\ch{^{\A}\X}};
        % Correction
        \foreach \i/\j/\X/\A in {7/7/C/11,4/7/C/14,7/8/N/12,6/8/N/13,3/8/N/16,
                                 2/8/N/17,1/9/O/19,5/9/O/15}
            \node[font=\footnotesize,inner sep=0.6pt,
                minimum size=5.5mm,
                rounded corners=1.4pt] at (m-\i-\j.center) {\ch{^{\A}\X}};
        \draw[thick,blue,shorten >=-2.9cm,opacity=0.7]
            (m-12-1.south west) -- (m-4-9.north east)
            node[rotate=45,below right] {$N=Z$}
            node[font=\small,rotate=45,above right,pos=1.04]
            {1$^{\text{ère}}$ bissectrice};
        \draw[very thick,red,->,shorten <=2.5mm,shorten >=2.5mm]
            (m-4-7.center) -- (m-5-8.center);
        \draw[very thick,red,->,shorten <=2.5mm,shorten >=2.5mm]
            (m-7-8.center) -- (m-6-7.center);
    \end{tikzpicture}
\end{figure}
\end{minipage}



